%% kapa: exam=1 round=19 year=2008 period=1 form=A subject=civil_law subject_ko=민법 seq=3 number=3
%% source: Q-Net 감정평가사 기출문제 (cst00309) artlSeq=1001962 file=extracted/1008/1001962/2008년 제19회 감정평가사 1차시험 기출문제/1교시.hwp
%% conversion: kordoc 4.24.51 -> tools/gen1.py (auto); review=proofread-sonnet; flags=none
\begin{question}{3}
의사표시에 관한 설명으로 옳은 것은? (다툼이 있는 경우에는 판례에 의함)
\begin{choices}
  \item 진의 아닌 의사표시에서 진의는 표의자가 진정으로 마음 속에서 바라는 사항을 의미한다.
  \item 통정허위표시의 무효로 대항할 수 없는 제3자가 되기 위해서는 선의이면 족하고, 무과실까지 요구되는 것은 아니다.
  \item 증여계약을 매매계약으로 가장한 경우에 매매계약뿐만 아니라 증여계약도 통정허위표시에 해당하여 무효이다.
  \item 의사표시의 상대방이 의사표시의 수령 당시에 미성년자인 경우, 표의자는 미성년자에 대하여 그 의사표시로 대항할 수 있다.
  \item 무효의 대상이 되는 강박행위는 불법적 해악의 고지로 상대방이 공포를 느끼는 정도에 이르는 것을 의미한다.
\end{choices}
\end{question}
\answer{2}
